\documentclass[12pt,a4paper]{article}

\usepackage[T2A]{fontenc}
\usepackage[utf8]{inputenc}
\usepackage[russian]{babel}
\usepackage{amsmath,amssymb,amsthm,mathtools}
\usepackage{geometry}
\usepackage{microtype}
\usepackage{enumitem}
\usepackage{xcolor}
\usepackage{tikz}
\usepackage{hyperref}

\geometry{left=25mm,right=20mm,top=22mm,bottom=24mm}
\hypersetup{
  colorlinks=true,
  linkcolor=blue!45!black,
  urlcolor=blue!55!black
}
\setlist{nosep}
\setlength{\parindent}{1.25em}
\setlength{\parskip}{0.15em}

\newtheorem{theorem}{Теорема}
\newtheorem{proposition}{Утверждение}
\theoremstyle{definition}
\newtheorem{definition}{Определение}
\theoremstyle{remark}
\newtheorem*{remark}{Замечание}

\DeclareMathOperator{\Sp}{Sp}
\DeclareMathOperator{\diag}{diag}
\DeclareMathOperator{\cond}{cond}
\newcommand{\CC}{\mathbb{C}}
\newcommand{\RR}{\mathbb{R}}
\newcommand{\ZZ}{\mathbb{Z}}
\newcommand{\eps}{\varepsilon}
\newcommand{\norm}[1]{\left\lVert #1\right\rVert}

\DeclareMathOperator{\ind}{ind}
\newcommand{\Conv}{\mathcal C}

\title{Операторы свёртки, алгебра Винера\\и индекс функции}
\author{Конспект лекции}
\date{10 октября 2026 года}

\begin{document}

\maketitle

\section{Оператор свёртки и его норма}

На предыдущей лекции было введено пространство
\[
  \ell^1(\ZZ)=\left\{a=(a_k)_{k\in\ZZ}:\sum_{k\in\ZZ}|a_k|<\infty\right\},
  \qquad \norm{a}_1=\sum_{k\in\ZZ}|a_k|.
\]
Для $a,x\in\ell^1(\ZZ)$ свёртка определяется формулой
\[
  (a*x)_k=\sum_{s\in\ZZ}a_{k-s}x_s.
\]
Она принадлежит $\ell^1(\ZZ)$ и удовлетворяет оценке
\[
  \norm{a*x}_1\leqslant\norm{a}_1\norm{x}_1.
\]
Свёртка коммутативна, ассоциативна и билинейна.

\begin{definition}
При фиксированном $a\in\ell^1(\ZZ)$ отображение
\[
  A_a\colon\ell^1(\ZZ)\longrightarrow\ell^1(\ZZ),\qquad A_ax=a*x,
\]
называется \emph{оператором свёртки}, а последовательность $a$ ---
его \emph{ядром}.
\end{definition}

Линейность свёртки по переменной $x$ означает линейность оператора:
\[
  A_a(\alpha x+\beta y)=\alpha A_ax+\beta A_ay.
\]
Доказанная оценка показывает, что этот оператор ограничен.

\begin{theorem}
Норма оператора свёртки на $\ell^1(\ZZ)$ равна норме его ядра:
\[
  \boxed{\norm{A_a}_{\ell^1\to\ell^1}=\norm{a}_1}.
\]
\end{theorem}

\begin{proof}
Из оценки свёртки для каждого $x$ имеем
\[
  \norm{A_ax}_1\leqslant\norm{a}_1\norm{x}_1.
\]
Следовательно, $\norm{A_a}\leqslant\norm{a}_1$.

Для обратной оценки возьмём единичную последовательность
\[
  e_k=\begin{cases}1,&k=0,\\0,&k\ne0.\end{cases}
\]
Она удовлетворяет равенствам $\norm{e}_1=1$ и $a*e=a$. Поэтому
\[
  \norm{a}_1=\norm{A_ae}_1
  \leqslant\norm{A_a}\norm{e}_1=\norm{A_a}.
\]
Объединяя обе оценки, получаем равенство норм.
\end{proof}

В частности, ядро однозначно восстанавливается по оператору:
\[
  \boxed{a=A_ae}.
\]
Равенство нормы здесь относится именно к действию на $\ell^1(\ZZ)$;
при смене рабочего пространства формулу для операторной нормы
нужно рассматривать заново.

\section{Алгебра операторов свёртки}

Обозначим через $\Conv$ множество всех операторов свёртки с ядрами
из $\ell^1(\ZZ)$. Умножением операторов считаем композицию:
$ABx=A(Bx)$.

\begin{proposition}
Для $a,b\in\ell^1(\ZZ)$ и $\alpha\in\CC$ справедливы равенства
\[
  A_a+A_b=A_{a+b},\qquad
  \alpha A_a=A_{\alpha a},\qquad
  A_aA_b=A_{a*b}.
\]
\end{proposition}

\begin{proof}
Первые два равенства следуют из билинейности:
\[
  (A_a+A_b)x=a*x+b*x=(a+b)*x,
  \qquad (\alpha A_a)x=(\alpha a)*x.
\]
Для произведения используем ассоциативность свёртки:
\[
  (A_aA_b)x=A_a(b*x)=a*(b*x)=(a*b)*x=A_{a*b}x.
\]
Так как равенства выполнены для всех $x$, соответствующие операторы
совпадают.
\end{proof}

Поскольку $a*b=b*a$, операторы свёртки коммутируют:
\[
  A_aA_b=A_bA_a.
\]
Единичный оператор также является оператором свёртки: $A_e=I$.

\begin{remark}
Ассоциативность композиции $(AB)C=A(BC)$ справедлива для любых
совместимых операторов. Коммутативность $AB=BA$ --- дополнительное
свойство рассматриваемых операторов свёртки; для произвольных
линейных операторов она неверна.
\end{remark}

\subsection{Изометрический изоморфизм с алгеброй ядер}

Рассмотрим отображение
\[
  \Phi\colon\ell^1(\ZZ)\longrightarrow\Conv,\qquad\Phi(a)=A_a.
\]
Оно обладает следующими свойствами:
\begin{enumerate}
  \item \emph{Биективность.} Каждый оператор из $\Conv$ имеет ядро
        по определению. Если $A_a=A_b$, то применением к $e$ получаем
        $a=b$.
  \item \emph{Линейность.}
        $\Phi(a+b)=\Phi(a)+\Phi(b)$ и
        $\Phi(\alpha a)=\alpha\Phi(a)$.
  \item \emph{Сохранение умножения.}
        $\Phi(a*b)=\Phi(a)\Phi(b)$.
  \item \emph{Сохранение нормы.}
        $\norm{\Phi(a)}=\norm{a}_1$.
\end{enumerate}
Таким образом, $\Phi$ --- изометрический изоморфизм алгебр.
В рукописи это отображение обозначено буквой $F$; здесь используется
$\Phi$, чтобы не смешивать его с матрицей ДПФ предыдущих лекций.

\begin{proposition}
С операторной нормой множество $\Conv$ является коммутативной
банаховой алгеброй с единицей.
\end{proposition}

\begin{proof}
Алгебраические свойства уже установлены. Проверим полноту.
Если $(A_{a^{(m)}})$ --- фундаментальная последовательность в
операторной норме, то
\[
  \norm{a^{(m)}-a^{(j)}}_1
  =\norm{A_{a^{(m)}}-A_{a^{(j)}}}\longrightarrow0.
\]
По полноте $\ell^1(\ZZ)$ существует $a\in\ell^1(\ZZ)$, для которого
$\norm{a^{(m)}-a}_1\to0$. Снова используя равенство норм, получаем
$A_{a^{(m)}}\to A_a$ в операторной норме. Предел принадлежит $\Conv$.
Наконец,
\[
  \norm{A_aA_b}=\norm{a*b}_1\leqslant\norm{a}_1\norm{b}_1
  =\norm{A_a}\norm{A_b}.
\]
\end{proof}

\section{Алгебра Винера}

\subsection{Абсолютно сходящиеся ряды на окружности}

Пусть
\[
  \Gamma=\{t\in\CC:|t|=1\}
\]
--- единичная окружность. Её можно параметризовать как $t=e^{i\varphi}$,
где $0\leqslant\varphi\leqslant2\pi$.

\begin{definition}
\emph{Алгеброй Винера} $W$ называется множество функций на $\Gamma$,
представимых в виде
\[
  a(t)=\sum_{k\in\ZZ}a_kt^k,\qquad
  \sum_{k\in\ZZ}|a_k|<\infty.
\]
\end{definition}

При $t=e^{i\varphi}$ это комплексный ряд Фурье
\[
  a(e^{i\varphi})=\sum_{k\in\ZZ}a_ke^{ik\varphi}.
\]
Отрицательные степени допустимы: на окружности $t$ никогда не равен нулю.

\begin{proposition}
Ряд, задающий элемент $W$, сходится абсолютно и равномерно на $\Gamma$.
Его сумма непрерывна.
\end{proposition}

\begin{proof}
Так как $|t|=1$, для каждого целого $k$ выполнено
\[
  |a_kt^k|=|a_k|.
\]
Числовой ряд $\sum_k|a_k|$ сходится. По признаку Вейерштрасса
функциональный ряд сходится равномерно и абсолютно. Каждый его член
непрерывен на $\Gamma$, поэтому непрерывна и сумма.
\end{proof}

\subsection{Связь с тригонометрической формой ряда Фурье}

Вещественную тригонометрическую запись обычно представляют в виде
\[
  \frac{A_0}{2}+\sum_{k=1}^{\infty}
    \bigl(A_k\cos(k\varphi)+B_k\sin(k\varphi)\bigr).
\]
Используя
\[
  \cos(k\varphi)=\frac{e^{ik\varphi}+e^{-ik\varphi}}2,
  \qquad
  \sin(k\varphi)=\frac{e^{ik\varphi}-e^{-ik\varphi}}{2i},
\]
получаем комплексные коэффициенты
\[
  c_0=\frac{A_0}{2},\qquad
  c_k=\frac{A_k-iB_k}{2},\qquad
  c_{-k}=\frac{A_k+iB_k}{2}\quad(k\geqslant1).
\]
Если $A_k,B_k$ вещественны, то $c_{-k}=\overline{c_k}$.

В рукописи переход также разбирается через вещественную и мнимую
части каждого коэффициента. Если $c_k=\alpha_k+i\beta_k$, то
\begin{align*}
  c_ke^{ik\varphi}
  &=(\alpha_k+i\beta_k)(\cos(k\varphi)+i\sin(k\varphi))\\
  &=\alpha_k\cos(k\varphi)-\beta_k\sin(k\varphi)\\
  &\quad+i\bigl(\beta_k\cos(k\varphi)+\alpha_k\sin(k\varphi)\bigr).
\end{align*}
При абсолютной сходимости эти равенства можно суммировать по всем
целым $k$.

\subsection{Однозначность коэффициентов и норма}

Коэффициенты представления функции в $W$ определяются однозначно.
Действительно, равномерная сходимость позволяет интегрировать ряд
почленно:
\[
  \boxed{a_m=\frac1{2\pi}\int_0^{2\pi}
      a(e^{i\varphi})e^{-im\varphi}\,d\varphi}.
\]
Здесь используется ортогональность экспонент:
\[
  \frac1{2\pi}\int_0^{2\pi}e^{i(k-m)\varphi}\,d\varphi
  =\begin{cases}1,&k=m,\\0,&k\ne m.\end{cases}
\]
Это пояснение обеспечивает корректность следующего определения нормы.

\begin{definition}
Нормой функции $a\in W$ называется сумма модулей её коэффициентов:
\[
  \boxed{\norm{a}_W=\sum_{k\in\ZZ}|a_k|}.
\]
\end{definition}

Норму функции в $W$ следует отличать от равномерной нормы на окружности:
\[
  \max_{t\in\Gamma}|a(t)|\leqslant\norm{a}_W.
\]
Равенство не обязательно. Например, для $a(t)=1+t-t^2$ норма Винера
равна $3$, тогда как
\[
  a(e^{i\varphi})=e^{i\varphi}(1-2i\sin\varphi),\qquad
  \max_\varphi|a(e^{i\varphi})|=\sqrt5.
\]
Этот дополнительный пример помогает различать два способа измерять
размер одной и той же функции.

\section{Операции в алгебре Винера}

Сложение и умножение функций в $W$ понимаются поточечно. При сложении
и умножении на скаляр коэффициенты преобразуются обычным образом:
\[
  (a+b)(t)=\sum_{k\in\ZZ}(a_k+b_k)t^k,\qquad
  (\alpha a)(t)=\sum_{k\in\ZZ}\alpha a_kt^k.
\]
Отсюда следуют оценки
\[
  \norm{a+b}_W\leqslant\norm{a}_W+\norm{b}_W,\qquad
  \norm{\alpha a}_W=|\alpha|\norm{a}_W.
\]

\begin{theorem}
Если $a,b\in W$, то $ab\in W$. Коэффициенты произведения равны свёртке
коэффициентов сомножителей:
\[
  (ab)(t)=\sum_{s\in\ZZ}c_st^s,\qquad
  \boxed{c_s=\sum_{j\in\ZZ}a_{s-j}b_j}.
\]
При этом $\norm{ab}_W\leqslant\norm{a}_W\norm{b}_W$.
\end{theorem}

\begin{proof}
Рассмотрим произведение рядов
\[
  a(t)b(t)=\left(\sum_{k\in\ZZ}a_kt^k\right)
          \left(\sum_{j\in\ZZ}b_jt^j\right).
\]
На окружности $|t|=1$ сумма модулей всех членов двойного ряда конечна:
\[
  \sum_{k\in\ZZ}\sum_{j\in\ZZ}|a_kb_jt^{k+j}|
  =\left(\sum_k|a_k|\right)\left(\sum_j|b_j|\right)<\infty.
\]
Поэтому члены можно группировать по одинаковой степени $s=k+j$:
\begin{align*}
  a(t)b(t)
  &=\sum_{k\in\ZZ}\sum_{j\in\ZZ}a_kb_jt^{k+j}\\
  &=\sum_{s\in\ZZ}\left(\sum_{j\in\ZZ}a_{s-j}b_j\right)t^s.
\end{align*}
Последовательность $(c_s)$ есть свёртка двух элементов $\ell^1(\ZZ)$.
По оценке, доказанной на предыдущей лекции,
\[
  \sum_s|c_s|\leqslant\left(\sum_k|a_k|\right)
                           \left(\sum_j|b_j|\right).
\]
Следовательно, $ab\in W$ и выполнена оценка нормы произведения.
\end{proof}

Единицей является постоянная функция $1$. Её последовательность
коэффициентов --- $e$, у которой $e_0=1$, а остальные координаты нулевые.
Коммутативность и ассоциативность умножения следуют из обычных свойств
поточечного произведения функций.

\subsection{Отображение функции в последовательность коэффициентов}

Определим
\[
  G\colon W\longrightarrow\ell^1(\ZZ),\qquad
  G(a)=(a_k)_{k\in\ZZ}.
\]
Его корректность следует из однозначности коэффициентов.
Инъективность следует из того, что одинаковые коэффициенты задают
одну и ту же функцию.
Сюръективность следует из определения $W$: любая последовательность
из $\ell^1$ задаёт абсолютно и равномерно сходящийся ряд на $\Gamma$.
Кроме того,
\begin{align*}
  G(a+b)&=G(a)+G(b),&G(\alpha a)&=\alpha G(a),\\
  G(ab)&=G(a)*G(b),&\norm{G(a)}_1&=\norm{a}_W.
\end{align*}
Таким образом, $G$ --- изометрический изоморфизм алгебр.

Полнота $W$ теперь получается из полноты $\ell^1(\ZZ)$. Действительно,
если $(a^{(m)})$ фундаментальна в норме $W$, то её последовательности
коэффициентов $G(a^{(m)})$ фундаментальны в $\ell^1$. Пусть их предел
равен $b$. Ряд с коэффициентами $b$ задаёт функцию $a\in W$, и
\[
  \norm{a^{(m)}-a}_W=\norm{G(a^{(m)})-b}_1\longrightarrow0.
\]
Итак, $W$ является коммутативной банаховой алгеброй с единицей.

\section{Символ оператора свёртки}

\begin{definition}
Если $A_a$ --- оператор свёртки с ядром $a=(a_k)_{k\in\ZZ}$, то функция
\[
  a(t)=\sum_{k\in\ZZ}a_kt^k,\qquad t\in\Gamma,
\]
называется \emph{символом оператора свёртки}.
\end{definition}

По определению символ принадлежит $W$, а потому непрерывен на
единичной окружности. Ядро, оператор и символ описывают один объект
в трёх изометрически изоморфных алгебрах:
\[
  \boxed{W\simeq\ell^1(\ZZ)\simeq\Conv}.
\]
Это соответствие сохраняет линейные операции, умножение и норму.

\begin{figure}[ht]
  \centering
  \begin{tikzpicture}[>=stealth]
    \node (w) at (0,0) {$a(t)\in W$};
    \node (l) at (4,0) {$(a_k)\in\ell^1(\ZZ)$};
    \node (c) at (8,0) {$A_a\in\Conv$};
    \draw[<->,thick,blue!60!black] (w) -- node[above] {$G$} (l);
    \draw[<->,thick,blue!60!black] (l) -- node[above] {$\Phi$} (c);
    \node[align=center] at (0,-0.9) {поточечное\\произведение};
    \node[align=center] at (4,-0.9) {свёртка\\последовательностей};
    \node[align=center] at (8,-0.9) {композиция\\операторов};
  \end{tikzpicture}
  \caption{Три представления и соответствующие им операции умножения}
\end{figure}

В частности, если операторам $A_a,A_b$ отвечают символы $a(t),b(t)$,
то сумма операторов имеет символ $a(t)+b(t)$, а композиция --- символ
$a(t)b(t)$. Равенство норм принимает вид
\[
  \norm{A_a}_{\ell^1\to\ell^1}
  =\norm{(a_k)}_1=\norm{a(\cdot)}_W.
\]
По символу можно восстановить коэффициенты интегральной формулой,
а по коэффициентам --- оператор свёртки.

\section{Индекс непрерывной функции на окружности}

\subsection{Определение через приращение аргумента}

Рассмотрим непрерывную функцию $a\colon\Gamma\to\CC$, которая не
обращается в нуль:
\[
  a(t)\ne0\qquad\text{для всех }t\in\Gamma.
\]
Параметризуем окружность положительно, то есть против часовой стрелки:
\[
  t=e^{i\varphi},\qquad0\leqslant\varphi\leqslant2\pi.
\]
Тогда $g(\varphi)=a(e^{i\varphi})$ задаёт замкнутую кривую, не
проходящую через начало координат. Замкнутость следует из
$g(0)=g(2\pi)=a(1)$.

Вдоль отрезка параметра можно выбрать непрерывный аргумент
$\theta(\varphi)$ так, чтобы
\[
  g(\varphi)=|g(\varphi)|e^{i\theta(\varphi)}.
\]
Существование такого выбора следует из непрерывности кривой и отсутствия
нулей: локальные ветви аргумента последовательно согласуются на
перекрывающихся участках. При этом $\theta(2\pi)$ не обязано совпадать
с $\theta(0)$.

\begin{definition}
\emph{Индексом} непрерывной ненулевой функции $a$ на $\Gamma$ называется
число
\[
  \boxed{\ind a=\frac{\theta(2\pi)-\theta(0)}{2\pi}}.
\]
Иначе говоря, это полное приращение аргумента при одном положительном
обходе окружности, делённое на $2\pi$.
\end{definition}

\begin{proposition}
Индекс является целым числом и не зависит от выбора непрерывного
аргумента вдоль кривой.
\end{proposition}

\begin{proof}
Так как $g(2\pi)=g(0)\ne0$, имеем
\[
  e^{i\theta(2\pi)}=e^{i\theta(0)}.
\]
Следовательно, разность $\theta(2\pi)-\theta(0)$ кратна $2\pi$.

Если $\widetilde\theta$ --- другой непрерывный аргумент, то для каждого
$\varphi$ разность $\widetilde\theta(\varphi)-\theta(\varphi)$ лежит
в $2\pi\ZZ$. Непрерывная функция на связном отрезке со значениями
в этом дискретном множестве постоянна. Поэтому приращения аргументов
совпадают.
\end{proof}

\begin{remark}
В определении нельзя просто вычесть главные значения аргумента в
концах пути: значения функции в начале и конце совпадают, и такое
вычитание потеряет число оборотов. Нужно отслеживать непрерывный
аргумент на всём отрезке $[0,2\pi]$.
\end{remark}

\subsection{Геометрический смысл и примеры}

Индекс --- число оборотов кривой $a(\Gamma)$ вокруг начала координат
с учётом направления и кратности прохождения. Обход против часовой
стрелки даёт положительный вклад, по часовой стрелке --- отрицательный.
Самопересечения кривой допускаются; проходить через нуль она не должна.

\begin{enumerate}
  \item Если $a(t)=t$, то $g(\varphi)=e^{i\varphi}$ и можно взять
        $\theta(\varphi)=\varphi$. Поэтому $\ind a=1$.
  \item Если $a(t)=t^2$, то $\theta(\varphi)=2\varphi$,
        и $\ind a=2$: окружность образа проходится дважды.
  \item Если $a(t)=t^{-1}$, то $\theta(\varphi)=-\varphi$,
        и $\ind a=-1$.
  \item Если $a(t)=2+t$, то образ --- окружность радиуса $1$
        с центром в точке $2$. Она лежит в правой полуплоскости.
        Аргумент возвращается к начальному значению, поэтому $\ind a=0$.
\end{enumerate}

Обобщая первые три примера, для любого целого $m$ получаем
\[
  \ind(t^m)=m.
\]

\begin{figure}[ht]
  \centering
  \begin{tikzpicture}[>=stealth,scale=0.9]
    \begin{scope}
      \draw[->,gray!65] (-1.35,0) -- (1.4,0);
      \draw[->,gray!65] (0,-1.2) -- (0,1.3);
      \draw[thick,blue!60!black] (0,0) circle (0.85);
      \draw[->,thick,blue!60!black] (20:0.85)
        arc[start angle=20,end angle=75,radius=0.85];
      \node[below left] at (0,0) {$0$};
      \node[align=center] at (0,-1.65) {$a(t)=t^2$\\два обхода; $\ind a=2$};
    \end{scope}
    \begin{scope}[xshift=4cm]
      \draw[->,gray!65] (-1.35,0) -- (1.4,0);
      \draw[->,gray!65] (0,-1.2) -- (0,1.3);
      \draw[thick,blue!60!black] (0,0) circle (0.85);
      \draw[->,thick,blue!60!black] (20:0.85)
        arc[start angle=20,end angle=75,radius=0.85];
      \node[below left] at (0,0) {$0$};
      \node[align=center] at (0,-1.65) {$a(t)=t$\\один обход; $\ind a=1$};
    \end{scope}
    \begin{scope}[xshift=7.4cm,scale=0.8]
      \draw[->,gray!65] (-0.4,0) -- (3.4,0);
      \draw[->,gray!65] (0,-1.5) -- (0,1.6);
      \draw[thick,blue!60!black] (2,0) circle (1);
      \draw[->,thick,blue!60!black] (2.94,0.342)
        arc[start angle=20,end angle=75,radius=1];
      \node[below left] at (0,0) {$0$};
      \node[below] at (2,0) {$2$};
      \node[align=center] at (1.6,-2.06) {$a(t)=2+t$\\$\ind a=0$};
    \end{scope}
  \end{tikzpicture}
  \caption{Индекс зависит от обхода образа, а не только от его формы}
\end{figure}

В частности, первые два примера имеют одно и то же множество значений
на комплексной плоскости, но разные индексы. Нужно учитывать, как
кривая проходится при изменении параметра.

\subsection{Аналитическая формула при дополнительной гладкости}

Рукопись заканчивается заголовком <<Аналитическая формула индекса>>;
самой формулы на скане нет. В качестве справочного дополнения приведём
её для случая, когда $g(\varphi)=a(e^{i\varphi})$ кусочно непрерывно
дифференцируема и нигде не равна нулю:
\[
  \boxed{\ind a=\frac1{2\pi i}\int_0^{2\pi}
    \frac{g'(\varphi)}{g(\varphi)}\,d\varphi}.
\]
На участках гладкости запишем $g=\rho e^{i\theta}$, $\rho>0$. Тогда
\[
  \frac{g'}g=\frac{\rho'}\rho+i\theta'.
\]
После интегрирования вещественная часть даёт
$\log\rho(2\pi)-\log\rho(0)=0$, а мнимая --- полное приращение
аргумента. Получается исходное определение индекса.

Принадлежность $W$ гарантирует непрерывность, но сама по себе не
гарантирует дифференцируемость. Поэтому определение через непрерывный
аргумент имеет более широкую область применимости, чем эта интегральная
запись с обычной производной.

\end{document}
